\documentclass[a4paper,11pt]{article}
\usepackage[a4paper,margin=2cm]{geometry}
\usepackage[T1]{fontenc}
\usepackage[utf8]{inputenc}
\usepackage[ngerman,english]{babel}
\usepackage{lmodern,microtype,xcolor,parskip,enumitem}
\usepackage[most]{tcolorbox}
\usepackage{titlesec}
\usepackage[hidelinks]{hyperref}
\definecolor{HeaderBlueDark}{RGB}{70,110,170}
\definecolor{RowBlueLight}{RGB}{235,243,255}
\definecolor{RowBlueDark}{RGB}{220,233,250}
\definecolor{SoftGray}{RGB}{90,90,90}
\setlength{\parindent}{0pt}
\setlist[itemize]{leftmargin=1.4em,itemsep=0.35em,topsep=0.35em}
\titleformat{\section}[block]{\Large\bfseries\color{HeaderBlueDark}}{}{0pt}{}
\newtcolorbox{exbox}{enhanced,breakable,colback=RowBlueLight,colframe=RowBlueDark,boxrule=0.4pt,arc=1.5mm,left=1.2mm,right=1.2mm,top=1mm,bottom=1mm,title={Beispiele \textnormal{(Examples)}},fonttitle=\bfseries,colbacktitle=RowBlueDark,coltitle=black}
\NewDocumentEnvironment{grammarentry}{mm+b}{\begin{tcolorbox}[enhanced,breakable,colback=white,colframe=HeaderBlueDark,title={#1\hfill\normalfont\footnotesize #2},colbacktitle=HeaderBlueDark,coltitle=white,fonttitle=\bfseries]#3\end{tcolorbox}}{}
\newcommand{\GermanText}[1]{\textbf{Deutsch: }#1\par}
\newcommand{\EnglishText}[1]{{\small\color{SoftGray}\textbf{English: }#1\par}}
\newcommand{\ExampleItem}[2]{\item \textbf{#1}\\{\small\color{SoftGray}#2}}
\title{Die Modalpartikel \glqq man\grqq\\{\large\textnormal{The modal particle man}}}
\date{}
\begin{document}
\maketitle
\tableofcontents
\section{Grundbedeutung \textnormal{(Basic meaning)}}
\begin{grammarentry}{man}{Modalpartikel / Abtönungspartikel}
\GermanText{Man ist als Abtönungspartikel eine norddeutsche umgangssprachliche Form, vor allem in lockeren Aufforderungen.}
\EnglishText{Man as a discourse particle is a northern German colloquial form, especially in casual commands.}
\end{grammarentry}
\section{Satzarten und Stellung \textnormal{(Sentence types and word order)}}
\begin{grammarentry}{Gebrauch}{Modalpartikel / Abtönungspartikel}
\GermanText{Regional norddeutsch; in der neutralen Standardsprache und für ÖSD-B1 nicht notwendig.}
\EnglishText{Regional northern German; not needed in neutral standard German or for ÖSD B1.}
\end{grammarentry}
\GermanText{Als Partikel ist man unflektierbar und verlangt keinen Kasus. Im üblichen Modalpartikelgebrauch steht es im Mittelfeld, nicht selbstständig im Vorfeld. Kontext und Betonung bestimmen die Gesprächswirkung.}
\EnglishText{As a particle, man does not inflect or govern a case. In its usual modal use it stands in the middle field, not independently in the prefield. Context and stress determine the conversational effect.}
\section{Typische Beispiele \textnormal{(Typical examples)}}
\begin{exbox}\begin{itemize}
\ExampleItem{Lass man gut sein!}{Come on, let it be! (regional)}
\ExampleItem{Komm man rein!}{Come on in! (regional)}
\ExampleItem{Mach man langsam!}{Take it easy! (regional)}
\ExampleItem{Setz dich man hin!}{Go on and sit down! (regional)}
\end{itemize}\end{exbox}
\section{Abgrenzung \textnormal{(Distinguishing other roles)}}
\begin{grammarentry}{Nicht jede Verwendung ist modal}{Modalpartikel / Abtönungspartikel}
\GermanText{Das gewöhnliche Indefinitpronomen man hat die Funktion eines Subjekts: Man lernt jeden Tag. Dieses Pronomen ist keine Modalpartikel.}
\EnglishText{The normal indefinite pronoun man functions as a subject: People learn every day. That pronoun is not a discourse particle.}
\end{grammarentry}
\begin{exbox}\begin{itemize}
\ExampleItem{Man lernt durch Übung.}{People learn through practice.}
\ExampleItem{Komm man rein!}{Come on in! (northern colloquial particle)}
\end{itemize}\end{exbox}
\section{Lernhinweis \textnormal{(Learning tip)}}
\GermanText{Modalpartikeln beeinflussen die Haltung oder Erwartung der sprechenden Person. Englische Entsprechungen sind situationsabhängig und häufig keine einzelnen Wörter.}
\EnglishText{Modal particles affect the speaker's attitude or expectations. English equivalents depend on context and often require a phrase rather than one word.}
\section{Quellen \textnormal{(Sources)}}
\begin{itemize}
\item IDS Mannheim: \url{https://grammis.ids-mannheim.de/terminologie/2}
\item Duden: \url{https://www.duden.de/sprachwissen/sprachratgeber/Modalpartikeln}
\end{itemize}
\end{document}
